% Reference deck: Chapter 3, Section 3.2.
\section[Properties]{Properties of Determinants}
\sectioncard{3.2}{Properties of Determinants}

\begin{frame}{Determinants and Elimination}
Cofactor expansion becomes expensive as the matrix grows.
\pause

Gaussian elimination can create a triangular matrix quickly, but row operations may change the determinant.
\medskip

The central question is therefore:
\[
\boxed{\text{How does each elementary row operation change }\det A?}
\]
\end{frame}

\begin{frame}{Three Row Operations}
Let $B$ be obtained from a square matrix $A$ by one row operation.
\theorembox{Theorem 3}{The determinant changes according to the operation:}
\renewcommand{\arraystretch}{1.35}
\begin{center}
\begin{tabular}{>{\raggedright\arraybackslash}p{.48\linewidth}c}
\toprule
Operation from $A$ to $B$ & Relationship \\
\midrule
$R_i\leftarrow R_i+cR_j$, $i\ne j$ & $\det B=\det A$ \\
$R_i\leftrightarrow R_j$ & $\det B=-\det A$ \\
$R_i\leftarrow kR_i$ & $\det B=k\det A$ \\
\bottomrule
\end{tabular}
\end{center}
\end{frame}

\begin{frame}{Row Replacement Preserves the Determinant}
Let
\[
A=\begin{bmatrix}1&2\\3&4\end{bmatrix},
\qquad \det A=-2.
\]
Apply $R_2\leftarrow R_2-3R_1$:
\[
B=\begin{bmatrix}1&2\\0&-2\end{bmatrix},
\qquad \det B=-2.
\]
\pause
Row replacement is the workhorse for determinant computation because it creates zeros without changing the determinant.
\end{frame}

\begin{frame}{A Row Interchange Reverses the Sign}
Interchange the two rows:
\[
A=\begin{bmatrix}1&2\\3&4\end{bmatrix}
\quad\longrightarrow\quad
B=\begin{bmatrix}3&4\\1&2\end{bmatrix}.
\]
Then
\[
\det A=-2,
\qquad
\det B=2=-\det A.
\]
\pause
After $r$ row interchanges, the accumulated factor is
\[
(-1)^r.
\]
\end{frame}

\begin{frame}{Equal Rows Force a Zero Determinant}
Suppose $A$ has two equal rows. Interchanging those rows leaves the matrix unchanged.
\pause

The interchange rule also says the determinant changes sign:
\[
\det A=-\det A.
\]
Therefore
\[
2\det A=0
\qquad\Longrightarrow\qquad
\boxed{\det A=0}.
\]
\end{frame}

\begin{frame}{Scaling One Row Scales the Determinant}
Let
\[
A=\begin{bmatrix}1&2\\3&4\end{bmatrix},
\qquad \det A=-2.
\]
Apply $R_1\leftarrow5R_1$:
\[
B=\begin{bmatrix}5&10\\3&4\end{bmatrix}.
\]
Then
\[
\det B=5(4)-10(3)=-10=5\det A.
\]
\pause
Factoring $k$ out of one row divides the visible row by $k$ and places a factor $k$ outside the determinant.
\end{frame}

\begin{frame}{Two Immediate Zero Tests}
\begin{block}{A zero row}
Scaling a zero row by any $k\ne1$ leaves the matrix unchanged, so
\[
\det A=k\det A.
\]
Hence $\det A=0$.
\end{block}
\pause
\begin{block}{One row is a multiple of another}
Use a row replacement to turn that row into a zero row. Row replacement preserves the determinant, so again $\det A=0$.
\end{block}
\end{frame}

\begin{frame}{Determinant by Row Reduction}
Reduce $A$ to an upper triangular matrix $U$.
\pause
Then
\[
\det U=u_{11}u_{22}\cdots u_{nn}.
\]
To recover $\det A$, keep a ledger:
\begin{itemize}
  \item row replacements contribute no factor;
  \item each row interchange contributes $-1$;
  \item each row scaling by $k$ contributes $k$ from the old determinant to the new one.
\end{itemize}
\end{frame}

\begin{frame}{Elimination Example: Setup}
Compute
\[
\det A,
\qquad
A=\begin{bmatrix}
1&4&2\\
2&8&9\\
1&7&0
\end{bmatrix}.
\]
Use row replacements:
\[
R_2\leftarrow R_2-2R_1,
\qquad
R_3\leftarrow R_3-R_1.
\]
\pause
Because both operations are replacements, the determinant does not change.
\end{frame}

\begin{frame}{Elimination Example: Finish}
The replacements give
\[
\det A=
\begin{vmatrix}
1&4&2\\
0&0&5\\
0&3&-2
\end{vmatrix}.
\]
Interchange rows $2$ and $3$:
\[
U=\begin{bmatrix}
1&4&2\\
0&3&-2\\
0&0&5
\end{bmatrix},
\qquad
\det U=1(3)(5)=15.
\]
\pause
One interchange gives
\[
\boxed{\det A=-\det U=-15}.
\]
\end{frame}

\begin{frame}{Elimination Without Swaps or Scaling}
Let
\[
B=\begin{bmatrix}
2&1&0\\
4&3&1\\
-2&2&5
\end{bmatrix}.
\]
Use only row replacements:
\[
\begin{aligned}
R_2&\leftarrow R_2-2R_1,\\
R_3&\leftarrow R_3+R_1,\\
R_3&\leftarrow R_3-3R_2.
\end{aligned}
\]
\pause
This produces
\[
U=\begin{bmatrix}2&1&0\\0&1&1\\0&0&2\end{bmatrix},
\qquad
\boxed{\det B=2(1)(2)=4}.
\]
\end{frame}

\begin{frame}{Avoid Scaling When Possible}
Standard reduced row echelon form encourages us to scale every pivot to $1$.
For determinants, this creates extra bookkeeping.
\pause

Instead, stop at any convenient echelon form:
\[
\begin{bmatrix}
p_1&*&\cdots&*\\
0&p_2&\cdots&*\\
\vdots&\ddots&\ddots&\vdots\\
0&\cdots&0&p_n
\end{bmatrix}.
\]
\pause
If only row replacements and $r$ interchanges were used, then
\[
\boxed{\det A=(-1)^r p_1p_2\cdots p_n}.
\]
\end{frame}

\begin{frame}{Combining Elimination and Cofactor Expansion}
For
\[
A=\begin{bmatrix}
0&1&2&-1\\
2&5&7&3\\
0&3&6&2\\
2&5&4&-2
\end{bmatrix},
\]
first use $R_4\leftarrow R_4-R_2$ to create
\[
\begin{bmatrix}
0&1&2&-1\\
2&5&7&3\\
0&3&6&2\\
0&0&-3&-5
\end{bmatrix}.
\]
\pause
Now the first column has a single nonzero entry, so a cofactor expansion reduces the problem to a $3\times3$ determinant.
\end{frame}

\begin{frame}{Column Operations}
Column operations have the same effects as the corresponding row operations:
\begin{itemize}
  \item adding a multiple of one column to another preserves the determinant;
  \item interchanging two columns reverses the sign;
  \item multiplying one column by $k$ multiplies the determinant by $k$.
\end{itemize}
\pause
This follows from the transpose property:
\[
\boxed{\det(A^T)=\det A}.
\]
A column operation on $A$ is the corresponding row operation on $A^T$.
\end{frame}

\begin{frame}{Transpose Property: Example}
Let
\[
A=\begin{bmatrix}1&2&0\\3&4&5\\0&6&7\end{bmatrix},
\qquad
A^T=\begin{bmatrix}1&3&0\\2&4&6\\0&5&7\end{bmatrix}.
\]
Expanding $A$ down its first column and $A^T$ across its first row gives the same cofactors and the same value:
\[
\det A=\det A^T.
\]
\pause
Rows and columns therefore play symmetric roles in determinant calculations.
\end{frame}

\begin{frame}{Multiplicative Property}
\theorembox{Theorem 6}{If $A$ and $B$ are $n\times n$ matrices, then}
\[
\boxed{\det(AB)=(\det A)(\det B)}.
\]
\pause
For
\[
A=\begin{bmatrix}2&1\\3&2\end{bmatrix},
\qquad
B=\begin{bmatrix}4&3\\1&2\end{bmatrix},
\]
we have $\det A=1$, $\det B=5$, and
\[
AB=\begin{bmatrix}9&8\\14&13\end{bmatrix},
\qquad
\det(AB)=117-112=5.
\]
\end{frame}

\begin{frame}{Consequences of Multiplicativity}
For square matrices of the same order,
\[
\det(A_1A_2\cdots A_m)
=\det A_1\det A_2\cdots\det A_m.
\]
In particular,
\[
\det(A^k)=(\det A)^k
\qquad(k=1,2,\ldots).
\]
\pause
Also, because multiplying the entire $n\times n$ matrix by $c$ scales each of its $n$ rows,
\[
\boxed{\det(cA)=c^n\det A}.
\]
\end{frame}

\begin{frame}{The Determinant Is Not Additive}
\begin{alertblock}{Common misconception}
In general,
\[
\det(A+B)\ne\det A+\det B.
\]
\end{alertblock}
\pause
For $A=B=I_2$,
\[
\det(A+B)=\det(2I_2)=4,
\]
while
\[
\det A+\det B=1+1=2.
\]
\end{frame}

\begin{frame}{Linearity in One Row or Column}
The determinant is linear in one row when all other rows stay fixed.
For example,
\[
\det\begin{bmatrix}
\vect u+\vect v\\
\vect r_2\\
\vdots\\
\vect r_n
\end{bmatrix}
=
\det\begin{bmatrix}
\vect u\\
\vect r_2\\
\vdots\\
\vect r_n
\end{bmatrix}
+
\det\begin{bmatrix}
\vect v\\
\vect r_2\\
\vdots\\
\vect r_n
\end{bmatrix}.
\]
\pause
The same statement holds for one column. It does not imply additivity in the entire matrix.
\end{frame}

\begin{frame}{Practice: Track the Operations}
Suppose $A$ is reduced to an upper triangular matrix $U$ by
\begin{enumerate}
  \item two row replacements;
  \item one row interchange;
  \item multiplying one row by $3$.
\end{enumerate}
If $\det U=-18$, find $\det A$.
\pause

The net relationship is
\[
\det U=(-1)(3)\det A.
\]
Therefore
\[
\boxed{\det A=6}.
\]
\end{frame}

\begin{frame}{Practice: Use Structure First}
Without expanding, evaluate
\[
\det\begin{bmatrix}
1&2&3&4\\
0&5&6&7\\
0&0&-2&8\\
0&0&0&3
\end{bmatrix}.
\]
\pause
The matrix is upper triangular, so
\[
\boxed{\det A=1(5)(-2)(3)=-30}.
\]
\end{frame}

